\PassOptionsToPackage{table,xcdraw}{xcolor}
\documentclass[11pt, a4paper]{article}

\usepackage[T1]{fontenc}
\usepackage[utf8]{inputenc}
\usepackage{tgpagella}
\usepackage[spanish, es-noshorthands]{babel}
\usepackage{amsmath, amssymb}
\usepackage[most]{tcolorbox}
\usepackage{tabularx}
\usepackage{tikz}
\usetikzlibrary{shapes.geometric, arrows.meta, positioning, calc}
\usepackage[margin=2.5cm]{geometry}
\usepackage{fancyhdr}

\pagestyle{fancy}
\fancyhf{}
\setlength{\headheight}{16pt}
\lhead{\textbf{UCB Sede Tarija} | Ing. Mecatrónica}
\rhead{IMT-221 | Reconciliación y Microdefensa}
\renewcommand{\headrulewidth}{0.4pt}
\definecolor{ucbblue}{RGB}{0, 51, 102}

\begin{document}

\begin{tcolorbox}[colback=red!5, colframe=red!70!black, title=\textbf{Documento Docente: Reconciliación de Evidencias (Triage)}]
    La microdefensa NO es automática ni es un examen de recuperación rutinario. Solo se gatilla ante incongruencias insalvables entre la demostración individual y la grupal[cite: 9].
\end{tcolorbox}

\begin{center}
\resizebox{0.85\textwidth}{!}{
\begin{tikzpicture}[auto, thick, node distance=2cm and 1.5cm]
    \node[draw, fill=blue!10, rounded corners, align=center] (exam) {Examen EC2 \\ (Evidencia Individual)};
    \node[draw, fill=blue!10, rounded corners, align=center, right=of exam] (report) {Reporte Hito 2 \\ (Evidencia Práctica Grupal)};
    
    \node[draw, fill=yellow!20, rounded corners, align=center, below right=1cm and -1.5cm of exam] (compare) {Comparación de Trazabilidad};
    
    \node[draw, fill=green!10, rounded corners, align=center, below left=1cm and -2cm of compare] (cong) {Congruente / Parcial};
    \node[draw, fill=red!10, rounded corners, align=center, below right=1cm and -2cm of compare] (incong) {Incongruente \\ (Contradicción real)};
    
    \node[draw, fill=orange!20, rounded corners, align=center, below=1cm of incong] (micro) {Microdefensa \\ Condicional};
    
    \draw[->] (exam) -- (compare);
    \draw[->] (report) -- (compare);
    \draw[->] (compare) -- (cong);
    \draw[->] (compare) -- (incong);
    \draw[->] (incong) -- (micro);
\end{tikzpicture}
}
\end{center}

\section*{Criterios de Gatillo de Microdefensa (Triggers)[cite: 9]}
\textbf{SÍ es candidato a Microdefensa:} Evidencia grupal muy sólida contrastada con evidencia individual pésima en el mismo dominio. Autoría incierta. Resultados reportados físicamente inverosímiles (fabricados)[cite: 9]. \\
\textbf{NO es candidato a Microdefensa:} Puntuación baja en examen por sí sola (va a Recuperatorio directo). Reporte incompleto. Errores aritméticos menores. Estudiante pidiendo "una oportunidad oral"[cite: 9].

\section*{Registro Oficial de Decisión Docente[cite: 9]}
\textbf{Estudiante:} \hrulefill \\[0.2cm]
\textbf{Puntaje Examen Oficial EC2:} \rule{2cm}{0.15mm} \quad \textbf{Puntaje Autoevaluado:} \rule{2cm}{0.15mm} \\[0.2cm]
\textbf{Dominio Analítico más Débil:} \hrulefill \\[0.2cm]
\textbf{Estatus del Reporte Grupal:} $\square$ Fuerte \quad $\square$ Débil \quad $\square$ Incompleto \quad $\square$ Inverosímil \\[0.2cm]
\textbf{Relación de Evidencias:} $\square$ Congruente \quad $\square$ Parcial \quad $\square$ Incongruente\\[0.4cm]
\textbf{¿Microdefensa Requerida?}\quad $\square$ SÍ \quad $\square$ NO \\[0.2cm]
\textbf{Razón detallada para Microdefensa:} \hrulefill \\[0.2cm]
\textbf{Conflicto de evidencia específico a resolver oralmente:} \hrulefill \\[0.4cm]
\textbf{¿Recuperación (Recovery) Requerida en Semana 10?}\quad $\square$ SÍ \quad $\square$ NO \\[0.2cm]
\textbf{Dominios obligatorios de Recuperación:} \hrulefill

\end{document}
